% Options for packages loaded elsewhere
\PassOptionsToPackage{unicode}{hyperref}
\PassOptionsToPackage{hyphens}{url}
\PassOptionsToPackage{dvipsnames,svgnames,x11names}{xcolor}
\documentclass[
  ngerman,
  10pt,
  a4paper,
]{article}
\usepackage{xcolor}
\usepackage[margin=22mm]{geometry}
\usepackage{amsmath,amssymb}
\setcounter{secnumdepth}{-\maxdimen} % remove section numbering
\usepackage{iftex}
\ifPDFTeX
  \usepackage[T1]{fontenc}
  \usepackage[utf8]{inputenc}
  \usepackage{textcomp} % provide euro and other symbols
\else % if luatex or xetex
  \usepackage{unicode-math} % this also loads fontspec
  \defaultfontfeatures{Scale=MatchLowercase}
  \defaultfontfeatures[\rmfamily]{Ligatures=TeX,Scale=1}
\fi
\usepackage{lmodern}
\ifPDFTeX\else
  % xetex/luatex font selection
  \setmainfont[Path=/Users/stefanhamann/.cache/codex-runtimes/codex-primary-runtime/dependencies/native/libreoffice-headless/libreoffice/LibreOfficeDev.app/Contents/Resources/fonts/truetype/,Extension=.ttf,BoldFont=DejaVuSans-Bold,ItalicFont=DejaVuSans-Oblique,BoldItalicFont=DejaVuSans-BoldOblique]{DejaVuSans}
  \setmonofont[Path=/Users/stefanhamann/.cache/codex-runtimes/codex-primary-runtime/dependencies/native/libreoffice-headless/libreoffice/LibreOfficeDev.app/Contents/Resources/fonts/truetype/,Extension=.ttf,BoldFont=DejaVuSansMono-Bold,ItalicFont=DejaVuSansMono-Oblique,BoldItalicFont=DejaVuSansMono-BoldOblique]{DejaVuSansMono}
\fi
% Use upquote if available, for straight quotes in verbatim environments
\IfFileExists{upquote.sty}{\usepackage{upquote}}{}
\IfFileExists{microtype.sty}{% use microtype if available
  \usepackage[]{microtype}
  \UseMicrotypeSet[protrusion]{basicmath} % disable protrusion for tt fonts
}{}
\makeatletter
\@ifundefined{KOMAClassName}{% if non-KOMA class
  \IfFileExists{parskip.sty}{%
    \usepackage{parskip}
  }{% else
    \setlength{\parindent}{0pt}
    \setlength{\parskip}{6pt plus 2pt minus 1pt}}
}{% if KOMA class
  \KOMAoptions{parskip=half}}
\makeatother
\ifLuaTeX
\usepackage[bidi=basic,shorthands=off]{babel}
\else
\usepackage[bidi=default,shorthands=off]{babel}
\fi
\ifPDFTeX
\else
\babelfont{rm}[Path=/Users/stefanhamann/.cache/codex-runtimes/codex-primary-runtime/dependencies/native/libreoffice-headless/libreoffice/LibreOfficeDev.app/Contents/Resources/fonts/truetype/,Extension=.ttf,BoldFont=DejaVuSans-Bold,ItalicFont=DejaVuSans-Oblique,BoldItalicFont=DejaVuSans-BoldOblique]{DejaVuSans}
\fi
\ifLuaTeX
  \usepackage{selnolig} % disable illegal ligatures
\fi
\setlength{\emergencystretch}{3em} % prevent overfull lines
\providecommand{\tightlist}{%
  \setlength{\itemsep}{0pt}\setlength{\parskip}{0pt}}
\usepackage{fvextra}
\RecustomVerbatimEnvironment{verbatim}{Verbatim}{breaklines=true,fontsize=\footnotesize}
\usepackage{xurl}
\usepackage{seqsplit}
\usepackage{adjustbox}
\usepackage{ragged2e}
\usepackage{titlesec}
\titleformat{\section}{\large\bfseries\raggedright}{}{0em}{}
\titleformat{\subsection}{\normalsize\bfseries\raggedright}{}{0em}{}
\DeclareRobustCommand{\codeword}[1]{\texttt{\seqsplit{#1}}}
\AtBeginEnvironment{longtable}{\small\RaggedRight}
\usepackage{newunicodechar}
\newunicodechar{𝒞}{\ensuremath{\mathcal C}}
\newunicodechar{𝔊}{\ensuremath{\mathfrak G}}
\newunicodechar{𝔤}{\ensuremath{\mathfrak g}}
\newunicodechar{𝓗}{\ensuremath{\mathcal H}}
\newunicodechar{ℋ}{\ensuremath{\mathcal H}}
\newunicodechar{𝔄}{\ensuremath{\mathfrak A}}
\newunicodechar{𝟙}{\ensuremath{\mathbf 1}}
\newunicodechar{𝕀}{\ensuremath{I}}
\newunicodechar{𝕋}{\ensuremath{\mathbb T}}
\newunicodechar{𝟘}{\ensuremath{\mathbf 0}}
\newunicodechar{⊞}{\ensuremath{\boxplus}}
\newunicodechar{∘}{\ensuremath{\circ}}
\newunicodechar{⟦}{\ensuremath{[\![}}
\newunicodechar{⟧}{\ensuremath{]\!]}}
\setlength{\emergencystretch}{4em}
\setlength{\parindent}{0pt}
\setlength{\parskip}{0.45em}
\AtBeginDocument{\hypersetup{colorlinks=true,linkcolor=black,urlcolor=blue}\pdfstringdefDisableCommands{\def\codeword#1{#1}}\RaggedRight}
\usepackage{fancyhdr}
\pagestyle{fancy}
\fancyhf{}
\fancyhead[L]{\small TFPT / Universalraum - v1.6.10}
\fancyfoot[R]{\thepage}
\setlength{\headheight}{15pt}
% The preserved full history now has three-digit page numbers.
\makeatletter
\renewcommand{\@pnumwidth}{2.6em}
\renewcommand{\@tocrmarg}{3.6em}
\makeatother

\usepackage{bookmark}
\IfFileExists{xurl.sty}{\usepackage{xurl}}{} % add URL line breaks if available
\urlstyle{same}
\hypersetup{
  pdflang={de-DE},
  colorlinks=true,
  linkcolor={Maroon},
  filecolor={Maroon},
  citecolor={Blue},
  urlcolor={Blue},
  pdfcreator={LaTeX via pandoc}}

\author{}
\date{}

\begin{document}

\section{Die fehlende halbe Drehung - einfach
erklärt}\label{die-fehlende-halbe-drehung---einfach-erkluxe4rt}

TFPT / Universalraum, v1.6.10, 15. September 2026.

\subsection{Das Bild}\label{das-bild}

Stell dir vier Stationen im Kreis vor. Zwischen benachbarten Stationen
können zwei Fermionen gemeinsam in einen anderen Zustand übergehen und
zurück. Bislang haben wir oft nur das Verhalten dieser Paare angeschaut.

Ein Fermion kann bei einem vollständigen Umlauf ein Minuszeichen
bekommen. Bei zwei Fermionen zusammen wird daraus wieder ein Plus. Der
Paarschatten kann deshalb nicht erkennen, ob sich darunter ein
gewöhnlicher oder ein verdrehter Kreis befindet.

Genau diese zusätzliche Information steckt bereits in einem älteren
TFPT-Modell: Vier Vierteldrehungen ergeben für das Fermion ein Minus;
erst acht ergeben wieder den Anfang. Das haben wir nun mit der
tatsächlichen inneren Paarwechselwirkung verbunden.

\subsection{Was dadurch besser wird}\label{was-dadurch-besser-wird}

Ohne die Verdrehung bleibt im gewählten Viererzyklus ein ganzer
Fermionbereich unbeteiligt. Wie ein Zimmer, das an keine Leitung
angeschlossen ist. Unterschiedliche Belegungen dieses Zimmers kosten
keine zusätzliche Energie. Deshalb kann das ganze Modell so keinen
einzigen eindeutig bevorzugten Grundzustand haben.

Mit der Verdrehung verschwindet dieses freie Zimmer. Für ausreichend
schwache Kopplung lässt sich nun beweisen: \textbf{Es gibt genau einen
energetisch günstigsten Zustand.} Dafür muss man nicht alle Zustände
einzeln ausrechnen. Ladungserhaltung und eine positive Schranke
schließen ganze Gruppen von Alternativen auf einmal aus.

Das ist mehr als eine Näherungsdiagonalisierung: Der Beweis
berücksichtigt auch die beliebig vielen möglichen Bosonbesetzungen. Er
gilt jedoch nur in dem ausdrücklich angegebenen kleinen Kopplungsbereich
dieses Viererzyklus.

\subsection{Und kann darin etwas wirklich angeregt
werden?}\label{und-kann-darin-etwas-wirklich-angeregt-werden}

Ja, ebenfalls innerhalb eines bewiesenen kleinen Bereichs. Entfernt man
ein ursprüngliches Fermion mit einer bestimmten Clock-Phase aus genau
diesem Grundzustand, landet fast die gesamte Antwort auf einer einzelnen
isolierten Energielinie. Je nach Clock-Phase gibt es zwei verschiedene
niedrige Linien. Für den passend aufgelösten Modus sind mehr als 99.992
Prozent des Entnahmespektrums kontrolliert.

Ein einzelner Stationsmodus mischt diese Phasen und sieht im Allgemeinen
beide Linien. Die Zahl beschreibt ein mathematisches Modellergebnis,
keine Messung eines echten Teilchens und noch keine Standardmodellmasse.

\subsection{Wo die große Hoffnung
liegt}\label{wo-die-grouxdfe-hoffnung-liegt}

Vielleicht haben wir nicht eine riesige zusätzliche Maschine übersehen,
sondern beim Übergang zu den sichtbaren Paargrößen eine kleine,
entscheidende Information weggeworfen. In diesem Test wirkt dasselbe
Umlaufvorzeichen auf die Kopplung, die Auswahl des Grundzustands und die
Anregungen. Es verbindet also mehrere Baustellen tatsächlich.

\subsection{Wo die Grenze liegt}\label{wo-die-grenze-liegt}

Wir haben die vier Stationen und ihre gemeinsame Quellenregel für diesen
Test festgelegt. Noch nicht gezeigt ist, dass die ursprüngliche TFPT
genau diese Regel erzeugt. Auch die bewiesene Kopplung ist viel kleiner
als unser bisheriger Prüfwert. Und vier Stationen sind keine abgeleitete
dreidimensionale Welt.

\textbf{Die Theory of Everything ist deshalb nicht fertig.} Neu ist eine
vollständig durchgerechnete lokale Verbindung, die vorher fehlte:
verdrehte Quelle, eindeutiger Modellgrundzustand und geladene Antwort
gehören nun zum selben System.

\subsection{Was die beiden neuen Texte zusätzlich
klären}\label{was-die-beiden-neuen-texte-zusuxe4tzlich-kluxe4ren}

Ihr Vorschlag, das Ganze über alle möglichen Handlungsfolgen und ihre
Antworten zu beschreiben, ist sinnvoll: Statt getrennte Schatten zu
sammeln, beschreibt man den gemeinsam ausführbaren Prozess. Aber ein
vollständiges Buch aller Antworten ist noch keine einfache Regel, die
dieses Buch erzeugt.

Die Zeitidee lässt sich ebenfalls bildlich prüfen: Ein Foto des
ruhigsten Zustands zeigt noch nicht, wie schnell das System schwingt,
wenn man es anstößt. Zwei Systeme können im Ruhezustand gleich aussehen
und verschieden schwingen. Eine aus dem Zustand berechnete „innere Uhr``
kann helfen, muss aber mit der tatsächlichen Bewegung verglichen werden.
Man darf die Bewegung nicht zuerst in den Zustand hineinschreiben und
sie anschließend als Entdeckung wieder auslesen.

Sehr hilfreich ist die andere Entlastung: Vielleicht müssen wir nicht
jede mikroskopische Leitung eindeutig festlegen. Entscheidend könnte
sein, dass eine ganze begründete Familie beim Vergrößern dieselben
physikalischen Gesetze zeigt. Genau das ist noch zu berechnen. Ein
kleines fertiges Modell ist dafür ein Anfang, nicht bereits die große
Welt.

\subsection{Was jetzt den Unterschied machen
würde}\label{was-jetzt-den-unterschied-machen-wuxfcrde}

Die Fortsetzung hat hier noch eine einfache Regel gefunden. Wenn linke
und rechte Nachbarstation verschieden stark beitragen dürfen,
funktioniert auch ein unverdrehtes Gegenmodell. Der Clock allein
bestimmt das also nicht. Fordert man aber, dass dieselbe
Links-rechts-Spiegelung auch den Paarübergang unverändert lässt, müssen
beide Seiten gleich stark beitragen. In dieser kleinen Quellenfamilie
bleibt dann ohne freie Zimmer nur die verdrehte Variante.

TFPT kennt bereits eine verwandte geometrische Spiegelung. Der
wichtigste nächste Test ist jetzt, ob sie tatsächlich \textbf{dieselbe
Spiegelung der Fermionorte und Paarübergänge} liefert. Das würde zwei
bislang eingesetzte Modellentscheidungen durch eine gemeinsame einfache
Regel ersetzen. Gezeigt ist bisher die bedingte Folgerung, nicht diese
ursprüngliche Identifikation. Danach muss beim Vergrößern herauskommen,
ob die Anregungen sich wie relativistische Materie bewegen. Die richtige
Messlatte bleibt gemeinsame physische Dynamik, nicht die Anzahl
passender Zahlen.

\end{document}
